\chapter{Clientes SQL}

Um \textbf{cliente SQL} é o programa pelo qual \emph{você} conversa com o
servidor: ele abre a conexão, envia os comandos SQL e exibe o resultado em
tabela. O servidor (MySQL) e o cliente são programas separados --- podem,
inclusive, estar em máquinas diferentes.

\section{Panorama de clientes}

\begin{table}[H]
\centering
\caption{Clientes SQL mais comuns para MySQL}
\label{tab:clientes}
\footnotesize
\begin{tabularx}{\textwidth}{@{}lllX@{}}
\toprule
\textbf{Cliente} & \textbf{Tipo} & \textbf{Plataformas} & \textbf{Custo} \\
\midrule
\comando{mysql} (CLI) & Terminal & Win / Linux / macOS & Gratuito, vem com o servidor \\
MySQL Workbench       & Desktop  & Win / Linux / macOS & Gratuito (GPL) \\
DBeaver               & Desktop  & Win / Linux / macOS & \emph{Community} gratuito; PRO pago \\
HeidiSQL              & Desktop  & Windows             & Gratuito \\
phpMyAdmin            & \textbf{Web} & navegador       & Gratuito \\
TablePlus             & Desktop  & Win / Linux / macOS & Versão gratuita limitada; pago \\
DBVisualizer          & Desktop  & Win / Linux / macOS & Versão gratuita limitada; Pro pago \\
Navicat               & Desktop  & Win / Linux / macOS & \textbf{Pago} (apenas teste grátis) \\
\bottomrule
\end{tabularx}
\end{table}

\begin{aviso}[Correção de um erro comum]
Navicat, DBVisualizer e SQLPro Studio \textbf{não são gratuitos}: são produtos
comerciais, com versão de teste ou edição gratuita limitada. Não os recomende
como ``gratuitos''.
\end{aviso}

\section{phpMyAdmin}

O phpMyAdmin \textbf{não é um programa que se instala no sistema operacional}
como os outros da lista: é uma \textbf{aplicação web escrita em PHP} que roda
num servidor web e é acessada pelo \textbf{navegador}. É por isso que ele
aparece junto do Apache e do PHP em pacotes como o XAMPP.

Ele foi escolhido para esta capacitação por três motivos: não exige instalação
na máquina do participante (basta uma URL), mostra a estrutura do banco
visualmente enquanto se aprende, e tem uma aba \textbf{SQL} onde os comandos
são digitados de verdade.

\section{Conectando}
\label{sec:conectando}

Nesta capacitação \textbf{ninguém instala nada}. O instrutor roda o MySQL e o
phpMyAdmin na máquina dele e publica apenas o phpMyAdmin na internet por um
túnel. O acesso é por uma URL, no navegador.

\begin{table}[H]
\centering
\caption{Dados de acesso da turma}
\label{tab:acesso}
\footnotesize
\begin{tabularx}{\textwidth}{@{}lX@{}}
\toprule
\textbf{Campo} & \textbf{Valor} \\
\midrule
URL do phpMyAdmin & Ditada pelo instrutor no início da aula \\
Usuário           & \comando{aluno} \\
Senha             & \comando{aluno} \\
\bottomrule
\end{tabularx}
\end{table}

\begin{dica}[Tela de aviso do túnel]
Na primeira visita, o serviço de túnel mostra uma página de aviso antes do
site. Clique em \emph{``Visit Site''} uma vez e siga; ela não aparece de novo.
\end{dica}

\subsection{Cada participante tem o seu próprio banco}
\label{sec:banco-proprio}

O servidor é um só, compartilhado pela turma inteira. Se todos trabalhassem no
mesmo banco, o \comando{CREATE TABLE} do segundo participante já falharia com
\comando{ERROR 1050: Table already exists}, e um \comando{DELETE} de qualquer
um apagaria o dado de todos.

Por isso \textbf{cada participante cria o próprio banco}, chamado
\comando{liga\_} seguido do seu primeiro nome:

\begin{center}
\comando{liga\_ana} \quad \comando{liga\_bruno} \quad \comando{liga\_carla}
\quad \comando{liga\_diego}
\end{center}

\begin{aviso}[Regra do nome]
Apenas letras minúsculas, sem acento e sem espaço: \comando{liga\_joao}, nunca
\comando{liga\_João} nem \comando{liga\_joao silva}.

Onde este guia escrever \comando{liga\_SEUNOME}, use o seu. O login
\comando{aluno} só tem permissão em bancos que comecem com \comando{liga\_};
qualquer outro nome retorna \comando{ERROR 1044: Access denied}. Isso é
proposital: protege o banco de demonstração do instrutor (\comando{cap}) e os
dados internos do MySQL.
\end{aviso}

\subsection{Teste de sanidade}

Abra a aba \textbf{SQL} e execute:

\begin{sqlcode}
SELECT VERSION(), NOW();
\end{sqlcode}

Se voltou a versão do servidor e a data e hora, está tudo certo.
